\documentclass[12pt]{article}

\usepackage[dvips,letterpaper,margin=0.75in,bottom=0.5in]{geometry}
\usepackage{cite}
\usepackage{slashed}
\usepackage{graphicx}
\usepackage{amsmath}
\usepackage{amssymb}
\usepackage{braket}
\usepackage{pythonhighlight}
\usepackage{mathtools}
\begin{document}

\newcommand{\ihbar}{\ensuremath{i \hbar}}
\newcommand{\Pss}{\ensuremath{\Psi^*}}
\newcommand{\dPsidt}{\ensuremath{ \frac{\partial \Psi}{\partial t} }}
\newcommand{\dPsidx}{\ensuremath{ \frac{\partial \Psi}{\partial x} }}
\newcommand{\ddPsidx}{\ensuremath{ \frac{\partial^2 \Psi}{\partial x^2} }}
\newcommand{\dPssdt}{\ensuremath{ \frac{\partial \Psi^*}{\partial t} }}
\newcommand{\dPssdx}{\ensuremath{ \frac{\partial \Psi^*}{\partial x} }}
\newcommand{\ddPssdx}{\ensuremath{ \frac{\partial^2 \Psi^*}{\partial x^2} }}

\newcommand{\dphidt}{\ensuremath{ \frac{d \phi}{dt} }}
\newcommand{\dpsidx}{\ensuremath{ \frac{d \psi}{dx} }}
\newcommand{\ddpsidx}{\ensuremath{ \frac{d^2 \psi}{dx^2} }}


\title{PHY 115L \\ Further Application \\ of Variational Methods }
\author{Michael Mulhearn}

\maketitle

\section{Notation}

We'll reserve $k$ for the wave number of a particle.  We will use $m$
and $n$ as dummy indices.  So we will use $M$ (capital) for the mass
of a particle to avoid confusion with the index $m$.  We'll use $N$
for the number of samples.  We'll use $a$ for the period of a
function.  Note that the numpy FFT library uses different notation,
taking $k$ as the index for Fourier coefficients (our $m$) and $n$ for
the number of samples (our $N$).  In our notation, the wave number
associated with the Fourier index $m$ is given by:
\begin{equation}
k = \frac{2 \pi m}{a}
\end{equation}

\section{TISE and System of Units}

\noindent
We will use a variational approach to find the ground state (and first
few exicted states)
$$\Psi_0(x), \; \Psi_1(x), \; \Psi_2(x), \; \ldots$$
of the Time-Independent
Schr\"odinger Equation (TISE):
$$\hat{H} \; \psi_0(x) = E_0 \; \psi_0(x)$$
where the Hamiltonian operator is defined by:
\begin{equation}
\hat{H} = -\frac{\hbar^2}{2 M} \; \frac{d^2}{dx^2} \, + \, V(x)
\end{equation}
We will consider potentials with a characteristic length scale $L$.
We will be considering {\bf non-periodic} wave function
(e.g. solutions to the harmonic oscillator).  To use the Fourier
series approach in this case, we take the period $a$ to be large
compared to $L$, and use the fact that the wave function approaches
zero to approximate a periodic function.

Based on our characteristic length scale $L$, a characteristic energy is:
\begin{equation}
\epsilon \equiv \frac{\hbar^2}{2 M L^2}
\end{equation}
Note well:  your program should not be calculating a numerical value for
$\epsilon$: that is the unit of energy for your program!  For example
$V(x) = 2$ in your program represents $V(x) = 2 \epsilon$
physically.  Using a system of units like this ensures that the
numbers we are calculating will be near $1$, which minimizes round-off
error.

With these units, we have:
\begin{equation}
\frac{\hat{H}}{\epsilon} = -L^2\frac{d^2}{dx^2} + \frac{V(x)}{\epsilon}
\end{equation}

\section{The Variational Method for Higher Energy States}

In the previous homework we found the ground state via:
$$\braket{\psi|\hat{H}|\psi} \geq E_0$$
for all $\psi$.  With knowledge that the ground state was even and the next excited state is odd, we can also determined the first excited state via:
$$\braket{\psi|\hat{H}|\psi} \geq E_1$$
for odd $\psi$.

We can generalize this to the following.  If we know that:
$$\braket{\Psi|m} = 0$$
for {\bf all} of the following:
$$m=0,1,\ldots,n$$
then we also have:  
$$\braket{\psi|\hat{H}|\psi} \geq E_{n+1}$$

We can, in principle at least, always arrange that our trial wave function is orthogonal to the previous discovered lower energy wave functions.  We use the projection operator:
$$P_m = \ket{m}\bra{m}$$
So, to remove the (previously discovered) ground state $\ket{\Psi_0}$ from
a trial wave function $\psi$ we calculate the component of $\ket{\Psi}$ parallel to $\ket{0}$, and subtract if from $\ket{\psi}$:
$$\ket{\Psi} \; \longrightarrow \; (1-P_m) \ket{\Psi} =
\ket{\Psi} - \braket{0|\Psi} \; \ket{0}$$
The inner product $\braket{0|\Psi}$ is easy to calculate in momentum space.  Suppose we have discovered $d_m$ so that:
$$\ket{0} = \sum_m d_m |p_m>$$
and have:
$$\ket{\Psi} = \sum_m c_m |p_m>$$
Then:
$$\braket{0|m} = \sum_m b_m^* c_m$$
And, for the case of real solutions with negative Fourier coefficients equal to the complex conjugate of the postive Fourier coefficients, we have:
$$\braket{0|m} = d_0^* c_0 + 2 \cdot \sum_{m>0} d_m^* c_m$$
Recall, that the python IRFFT library provides the Fourier coefficents with slightly different normalization:
\begin{equation}
A_m \equiv \frac{c_m}{\sqrt{a/N}} \exp\left(i \frac{2\pi m x_0}{a}\right)
\end{equation}
and
\begin{equation}
B_m \equiv \frac{d_m}{\sqrt{a/N}} \exp\left(i \frac{2\pi m x_0}{a}\right)
\end{equation}
so the dot product using the coefficients from IRFFT becomes:
$$\braket{0|m} = \Delta a \left\{ B_0^* A_0 + 2 \cdot \sum_{m>0} B_m^* A_m \right\}$$
where, as before,
$$\Delta a = \frac{a}{N}$$.
In principle, we can substract any number of lower energy states to
find the next highest energy states, but round off errors will
eventually dominate.  Note that when subtracting multiple lower energy
states, you only need to normalize after subtracting the last
component.  One can comine the even versus odd trick to make this more
efficient (e.g. subtracting the ground state from an even trial
function would lead to the 2nd excited state, if the energy
eigenstates alternate between even and odd.)

The numerical algorith proceeds much the same way, with one modification: after {\bf each} random variation of the wave function, you remove the lower energy components, renormalize, and then check for a drop an energy.

\section{The Hydrogen Atom}

We can apply our numerical technique to find the radial wave functions of
the Hydrogen Atom.  In this case, we'll take our characteristic length scale $L$ to be the Bohr radius:
$$L \equiv \frac{4 \pi \epsilon_0 \hbar^2}{m e^2}$$
As always, this quantity should not be computed in your computer program, it will simple be set to one in every equation where it appears.
We continue to have our energy scale defined by:
$$\epsilon = \frac{\hbar^2}{2mL^2}$$
but with $L$ now defined we have an alternative form:
$$\epsilon = \frac{m}{2\hbar^2} \left( \frac{e^2}{4\pi\epsilon_0}\right)^2
= 2 \epsilon L $$
Another useful quantity is:
$$\frac{e^2}{4\pi\epsilon_0} = 2 \epsilon L$$
Using this system of units, the protential for the Hydrogen atom becomes:
$$\frac{V(r)}{\epsilon} = - \frac{2L}{r}$$
The allowed energy levels are:
$$E_n =  \epsilon / n^2$$
where $n$ is the principle quantum number.  The additional quantum numbers $l$ and $m$ refer to the angular momentum of the state.

\section{Radial Solutions}

Using the method of separation of variables (e.g. Griffiths Chapter 4) we write the total wave function as:
$$\Psi_{nlm}(r,\theta,\phi) = \frac{u_nl(r)}{r} Y_{lm}(\theta,\phi)$$
for quantum numbers $n,l,m$.  The radial part satisfies 
\begin{equation*}
  \frac{\hat{E}}{\epsilon} u(r) = -L^2\frac{d^2 u}{dr^2} + \left(-\frac{2L}{r}+ l(l+1) \frac{L^2}{r^2}\right)
\end{equation*}





The first few radial solutions are:
$$u_{10} = \frac{2}{\sqrt{L}} \; \frac{r}{L} \; \exp\left(-\frac{r}{L}\right)$$
and:
$$u_{20} = \frac{2}{\sqrt{L}} \; \frac{r}{2L} \; \left(1 - \frac{r}{2L} \right)\; \exp\left(-\frac{r}{2L}\right)$$
and
$$u_{21} = \frac{2}{\sqrt{L}} \; \left( \frac{r}{2L} \right)^2 \; \exp\left(-\frac{r}{2L}\right)$$
Notice that none of these solutions are even or odd, but they do all vanish at $r=0$.

We can find these solutions using just a few modifications to our algorithm from the previous homework, via the replacement:
$$x \longrightarrow r$$
First, we must impose the boundary condition
$u(0)=0$
This is easy to implement in our Fourier expansion.  Because the sine vanishes at $x=0$, we will meet this condition provided:
$$c_0 = - 2 \cdot \sum_{m>0} {{\rm Re}\;{c_m}}$$
Recall that for a real function $c_0$ is real.  So we only need to
modify our algorithm to not include $c_0$ in our random variations.
We fix the imaginary part of $c_0$ to zero, but after each variation,
we modify the real part of $c_0$ by the equation above.

We'll also need to change $x_0$ from $-a/2$ to $0$, appropriate for the radius $r$.


\section{Homework Problems}

\noindent
{\bf Please submit one PDF file (including output) and one python
  notebook file (.ipynb) of your completed assignment.}  If using
C/C++, submit text or PDF of your output and the source code.


\noindent
{\bf Problem 1:} Building on your results from HW2, find the second and third excited state of the Harmonic oscillator potential.\\[3pt]

\noindent
{\bf Problem 2:} For the hydrogen, and using their known
analytic functions, plot $V(x)$, $u_{10}(r)$, $u_{20}(r)$, and $u_{21}(r)$ versus $r$ in our system of units. For better visiblity, offset the
$y$-values by the energy of each state.  So e.g. plot $u_{10}(x)=0$ at
$V(x)/\epsilon=-1$.

{\bf Problem 3:} Find the Fourier coefficients of $u_{10}$ and $u_{20}$



For the hydrogen, and using their known
analytic functions, plot $V(x)$, $u_{10}(r)$, $u_{20}(r)$, and $u_{21}(r)$ versus $r$ in our system of units. For better visiblity, offset the
$y$-values by the energy of each state.  So e.g. plot $u_{10}(x)=0$ at
$V(x)/\epsilon=-1$.


\noindent
{\bf Problem 3:} Use the variational method to calculate the ground state of the hydrogen atom.  Compare your result in terms of the calculated energy and 



\noindent
{\bf Problem 3:} By turning on the centrifugal term proportional to $l(l+1)$, find the energy for the first few excited states, corresponding to $u_{21}$, $u_{32}$, and $u_{43}$.

    


    



\end{document}
